\scanpage{41}
% archive foliation: 40
\begin{center}
UNIVERSITE DES SCIENCES\\
ET\\
TECHNIQUES DU LANGUEDOC
\end{center}
\hrule

Montpellier, le 9 mars 1982.

Service du Personnel Enseignant

MCS/FR.

N°1317.

Le Président

à

Mesdames et Messieurs\\
les enseignants de l'U.S.T.L.

\uline{OBJET} : Elections au Comité Technique Paritaire Central des Personnels Enseignants de statut universitaire -\\
Scrutin par correspondance de mars 1982.

\uline{Réf.} : Décret n°82.139 du 5 février 1982.\\
Arrêté du 8 février 1982.

J'ai l'honneur de vous adresser le matériel électoral nécessaire aux opérations de vote par correspondance, comprenant :

- 1 notice explicative ;

- 2 bulletins de vote (constitués par les listes des candidats)\\
\hspace*{2em}. SGEN - CFDT\\
\hspace*{2em}. SNESUP

- 2 professions de foi (1 pour chaque liste)

- jeu de 3 enveloppes

\hspace*{2em}. enveloppe n°1 blanche, ne comportant aucune marque ou distinction permettant d'en déterminer l'origine ;

\hspace*{2em}. enveloppe n°2 portant les mentions : nom, prénom, grade, affectation, signature ;

\hspace*{2em}. enveloppe n°3 libellée au nom de M. le Recteur, Président de la Commission de Dépouillement.

Je vous précise que les votes devront être adressés, \uline{par voie postale à M. le Recteur de l'Académie de Montpellier, la date limite de réception étant fixée au 30 mars 1982 à 17 heures,le cachet de la poste faisant foi}.

Les votes parvenus hors délais seront renvoyés aux électeurs.

Le dépouillement public se déroulera au Rectorat, Salle des Commissions, le \uline{31 mars 1982 à 14 heures}.

Le Président,

\unsure{L. Thaler}

\uline{L. THALER}.
